% 分野別類題: 逆行列・行列式・行列方程式 解答例
% コンパイル: TEXINPUTS="../../../00_common:$TEXINPUTS" latexmk -xelatex answers.tex
\documentclass[a4paper,11pt]{article}
\usepackage{preamble}
\usepackage{shoakubox}

\begin{document}

\kamokumei{類題演習 解答例 --- 逆行列・行列式・行列方程式}

\daimon{【解法】3 次正方行列の行列式・逆行列・行列方程式の解き方と導出}

\noindent
\textbf{(1) 行列式（余因子展開）}\quad
3 次正方行列 $A = (a_{ij})$ の行列式は、第 1 行に関する余因子展開により
\[
\det A = a_{11}A_{11} + a_{12}A_{12} + a_{13}A_{13}, \qquad
A_{ij} = (-1)^{i+j}M_{ij}
\]
と表される。ここで $M_{ij}$ は $A$ から第 $i$ 行・第 $j$ 列を除いた $2\times 2$ 小行列式（小行列式）であり、$A_{ij}$ を $(i,j)$ 余因子と呼ぶ。どの行・列に関して展開しても値は一致する。掃き出し法（行基本変形）で三角化しても、行の入れ替えで符号が変わる・行の定数倍で行列式もその定数倍になる、という規則の下で同じ値が得られる。

\noindent
\textbf{(2) 逆行列（余因子行列・掃き出し法）}\quad
$\det A \neq 0$ のとき $A$ は正則であり、逆行列は余因子行列（の転置、随伴行列）を用いて
\[
A^{-1} = \frac{1}{\det A}\,\operatorname{adj}A, \qquad
\operatorname{adj}A = (A_{ij})^{T}
\]
と与えられる。これは $A \cdot \operatorname{adj}A = (\det A) I$ という恒等式（各行と対応しない余因子の組み合わせによる余因子展開が $0$ になることから従う）の帰結である。実務上は掃き出し法で $(A \mid I) \to (I \mid A^{-1})$ と変形して求めてもよい。

\noindent
\textbf{(3) 行列方程式 $AX = B$（$XA = B$）}\quad
$A$ が正則ならば、両辺に左から $A^{-1}$ を掛けて
\[
AX = B \;\Longrightarrow\; X = A^{-1}B
\]
と一意に解ける（$XA = B$ の場合は右から掛けて $X = BA^{-1}$）。$A$ が正則でなければ、一般には解が存在しないか無数に存在する。

\noindent
\textbf{(4) パラメータを含む正則条件}\quad
パラメータ $t$ を含む行列 $A(t)$ が正則でない（逆行列を持たない）ための必要十分条件は $\det A(t) = 0$ である。したがって $\det A(t)$ を $t$ の多項式として計算し、その根を求めればよい。

\clearpage
\begin{mondaiboxS}{問1}
次の行列 $A$ について、行列式 $\det A$ と逆行列 $A^{-1}$ を求めよ。
\[
A =
\begin{pmatrix}
1 & 2 & 0 \\
0 & 1 & 3 \\
2 & 0 & 1
\end{pmatrix}
\]
\end{mondaiboxS}
\kaitou
第 1 行に関して余因子展開する。
\[
\det A = 1\cdot\begin{vmatrix}1 & 3\\ 0 & 1\end{vmatrix}
- 2\cdot\begin{vmatrix}0 & 3\\ 2 & 1\end{vmatrix}
+ 0\cdot\begin{vmatrix}0 & 1\\ 2 & 0\end{vmatrix}
= 1\cdot 1 - 2\cdot(-6) + 0 = 1 + 12 = 13
\]
よって
\[
\boxed{\; \det A = 13 \;}
\]
次に各余因子を計算する。
\[
A_{11}=1,\ A_{12}=6,\ A_{13}=-2,\quad
A_{21}=-2,\ A_{22}=1,\ A_{23}=4,\quad
A_{31}=6,\ A_{32}=-3,\ A_{33}=1
\]
（例: $A_{12} = -\begin{vmatrix}0&3\\2&1\end{vmatrix} = -(0-6)=6$、$A_{23} = -\begin{vmatrix}1&2\\2&0\end{vmatrix} = -(0-4)=4$。）
余因子行列を転置して随伴行列を作ると
\[
\operatorname{adj}A =
\begin{pmatrix}
1 & -2 & 6 \\
6 & 1 & -3 \\
-2 & 4 & 1
\end{pmatrix}
\]
したがって
\[
\boxed{\; A^{-1} = \frac{1}{13}
\begin{pmatrix}
1 & -2 & 6 \\
6 & 1 & -3 \\
-2 & 4 & 1
\end{pmatrix} \;}
\]
（検算: 例えば $A$ の第 1 行 $(1,2,0)$ と $A^{-1}$ の第 1 列 $(1,6,-2)/13$ の内積は $(1+12+0)/13=1$、第 2 列 $(-2,1,4)/13$ との内積は $(-2+2+0)/13=0$ となり、$AA^{-1}=I$ を満たす。他の成分も同様に確認できる。）

\clearpage
\begin{mondaiboxS}{問2}
次の行列方程式 $AX = B$ を満たす行列 $X$ を求めよ。
\[
A =
\begin{pmatrix}
1 & 1 & 0 \\
0 & 1 & 1 \\
1 & 0 & 1
\end{pmatrix},
\qquad
B =
\begin{pmatrix}
4 & 3 \\
3 & 5 \\
1 & 6
\end{pmatrix}
\]
\end{mondaiboxS}
\kaitou
まず $A$ が正則であることを確かめ、$A^{-1}$ を求める。第 1 行に関して余因子展開すると
\[
\det A = 1\cdot\begin{vmatrix}1&1\\0&1\end{vmatrix} - 1\cdot\begin{vmatrix}0&1\\1&1\end{vmatrix} + 0
= 1\cdot 1 - 1\cdot(-1) = 2 \neq 0
\]
であるから $A$ は正則である。余因子を計算すると
\[
A_{11}=1,\ A_{12}=1,\ A_{13}=-1,\quad
A_{21}=-1,\ A_{22}=1,\ A_{23}=1,\quad
A_{31}=1,\ A_{32}=-1,\ A_{33}=1
\]
より随伴行列は
\[
\operatorname{adj}A =
\begin{pmatrix}
1 & -1 & 1 \\
1 & 1 & -1 \\
-1 & 1 & 1
\end{pmatrix},
\qquad
A^{-1} = \frac{1}{2}
\begin{pmatrix}
1 & -1 & 1 \\
1 & 1 & -1 \\
-1 & 1 & 1
\end{pmatrix}
\]
$A$ は正則だから $X = A^{-1}B$ と一意に定まる。$A^{-1}$ の各行と $B$ の各列の内積を計算すると
\[
X = \frac{1}{2}
\begin{pmatrix}
1 & -1 & 1 \\
1 & 1 & -1 \\
-1 & 1 & 1
\end{pmatrix}
\begin{pmatrix}
4 & 3 \\
3 & 5 \\
1 & 6
\end{pmatrix}
= \frac{1}{2}
\begin{pmatrix}
4-3+1 & 3-5+6 \\
4+3-1 & 3+5-6 \\
-4+3+1 & -3+5+6
\end{pmatrix}
= \frac{1}{2}
\begin{pmatrix}
2 & 4 \\
6 & 2 \\
0 & 8
\end{pmatrix}
\]
したがって
\[
\boxed{\; X =
\begin{pmatrix}
1 & 2 \\
3 & 1 \\
0 & 4
\end{pmatrix} \;}
\]
（検算: $AX$ を直接計算すると、第 1 行 $(1,1,0)$ と $X$ の各列の内積は $(1+3,\,2+1)=(4,3)$、第 2 行 $(0,1,1)$ とは $(3+0,\,1+4)=(3,5)$、第 3 行 $(1,0,1)$ とは $(1+0,\,2+4)=(1,6)$ となり、これは $B$ に一致する。）

\clearpage
\begin{mondaiboxS}{問3}
パラメータ $t$ を含む次の行列 $A(t)$ について、$A(t)$ が正則でない（逆行列を持たない）ような $t$ の値をすべて求めよ。
\[
A(t) =
\begin{pmatrix}
t & 1 & 1 \\
1 & t & 1 \\
1 & 1 & t
\end{pmatrix}
\]
\end{mondaiboxS}
\kaitou
$A(t)$ が正則でないための必要十分条件は $\det A(t) = 0$ であるから、まず $\det A(t)$ を計算する。第 2 行・第 3 行からそれぞれ第 1 行を引いても行列式の値は変わらないから、
\[
\det A(t) =
\begin{vmatrix}
t & 1 & 1 \\
1-t & t-1 & 0 \\
1-t & 0 & t-1
\end{vmatrix}
\]
第 2 行・第 3 行から共通因子 $(t-1)$ をくくり出すと
\[
\det A(t) = (t-1)^{2}
\begin{vmatrix}
t & 1 & 1 \\
-1 & 1 & 0 \\
-1 & 0 & 1
\end{vmatrix}
\]
右辺の $3\times3$ 行列式を第 1 行で余因子展開すると
\[
\begin{vmatrix}
t & 1 & 1 \\
-1 & 1 & 0 \\
-1 & 0 & 1
\end{vmatrix}
= t(1\cdot1 - 0\cdot0) - 1\bigl((-1)\cdot1 - 0\cdot(-1)\bigr) + 1\bigl((-1)\cdot0 - 1\cdot(-1)\bigr)
= t + 1 + 1 = t+2
\]
したがって
\[
\det A(t) = (t-1)^{2}(t+2)
\]
（検算: $t=0$ を代入すると $\det A(0) = (-1)^2\cdot 2 = 2$。一方、直接展開すると $A(0)=\begin{pmatrix}0&1&1\\1&0&1\\1&1&0\end{pmatrix}$ で、$\det = 0(0-1)-1(0-1)+1(1-0)=0+1+1=2$ となり一致する。）

$A(t)$ が正則でないのは $\det A(t) = 0$ のときであり、$(t-1)^{2}(t+2)=0$ を解くと
\[
\boxed{\; t = 1 \ \text{（重解）}, \quad t = -2 \;}
\]
（検算: $t=1$ のとき $A(1)=\begin{pmatrix}1&1&1\\1&1&1\\1&1&1\end{pmatrix}$ は全行が等しく階数 $1$ で正則でない。$t=-2$ のとき $A(-2)=\begin{pmatrix}-2&1&1\\1&-2&1\\1&1&-2\end{pmatrix}$ は各行の和が $0$ となるため $(1,1,1)^{T}$ が核に属し、正則でない。）

\end{document}
